This thesis focuses on the use of reinforcement learning for adversarial planning within the framework of realistic guidance, navigation, and controls for orbital spacecraft. Traditional approaches based on classical optimization methods can become brittle when faced against intelligent and non-cooperative opponents which can have shifting strategies. Reinforcement learning algorithms can be useful tools to fill in this gap, so long as they're used within the right use case and context that mitigates these algorithms' lack of convergence guarantees, proneness to instability, non-stationarity, and limited utilization outside the training dataset.

With this in mind, this thesis proposes the Phased Adversarial Reinforcement Learning plus Behavior Cloning (PARL+BC, for short) framework to train robust policies in a two-player, zero-sum Target-Attacker-Defender (TAD) orbital defense scenario. This framework is built based on the well-known Proximal Policy Optimization (PPO) online reinforcement learning algorithm. We train attackers by freezing existing opponent policies and using randomized mixtures of these, with the intent of mitigating non-stationarity during learning and avoiding the instabilities that are often seen on simultaneous play methods. 

These learned policies train within a closed-loop aerospace guidance, navigation, and control architecture. The control piece comes in the form of a Control Barrier Function-Quadratic Program (CBF-QP) for velocity saturation within a predefined safety boundary. Moreover, the navigation piece, which is meant to showcase the feasibility of the framework under partial observability simulating real-life sensor limitations, is addressed by doing behavior cloning from observations obtained from an Extended Kalman Filter in the loop.

We trained three different test cases that resulted in multiple attacker and defender policies across consecutive matchups, which showed stability during learning and convergent behavior.

Furthermore, the policies are evaluated via Monte Carlo rollouts across full and partial observability, varying arena sizes, and the proper dispersions to randomize the starting states of each agent as well as their EKF parameters. While policies are trained under Hill-Clohessy-Wiltshire (HCW) linear time invariant relative orbital dynamics, we evaluated the policies in both HCW and elliptical Linear-Time-Varying (LTV) dynamics to see if the learned policies can transfer to time-varying dynamics that were not seen during training.

The Monte Carlo results show that our proposed framework produces coherent and competitive defender and attacker policies for the proposed zero-sum game. Within a small, 20 m training arena, full-state policies reached up to a 79.9\% defender win rate against a competitive RL attacker with full state access to the opponent's state, and 74.7\% in the EKF-based policies with bearings-only measurements of the opponent's state, a 6.51\% relative drop in performance relative to the fully observable policies.

When evaluating in a larger 100m radius arena, not used or seen by the policies during training, full-state results achieved up to a 66.2\% defender win rate, a 17.15\% relative drop in performance from the 20m training arena, whereas EKF-based policies struggled more with a 47.4\% defender win rate, a 28.4\% relative drop from the full state case. 

Evaluating all of the above across both HCW and elliptic LTV dynamics, even though the policies were never exposed to the latter during training, resulted in at most a 1.29\% relative change between results. Therefore, within the cases tested in this work, the policies trained under linear time-invariant dynamics transferred effectively to the linear time-varying dynamics used during evaluation.

Throughout the Monte Carlo runs, median policy inference time reached up to 0.145 ms per timestep on the full-state policies and 3.91 ms per timestep on the EKF-based policies, the latter of which includes the EKF compute time in the loop.

All in all, the results obtained via our proposed PARL+BC method and its evaluation show that reinforcement learning can be integrated with classical navigation and control architectures in the loop in constrained adversarial orbital scenarios, resulting in robust and computationally efficient policies that work under partial observability and can extrapolate beyond the observations obtained through the training dataset.
